\chapter{Benchmark design}\label{chap:map:benchmark}

\section{Methodology}

To ensure fair and accurate benchmarks, we followed a rigorous methodology. We measured the execution time of various operations (insertion, deletion, lookup) on our data structures using the \texttt{clock()} function.

For map benchmarks, we used random \texttt{uint64\_t} keys and values. To isolate the performance of the data structure itself from the hash function, we used the identity function as the hash function (assuming the keys are already random).

We performed benchmarks on:
\begin{itemize}
    \item \textbf{Unsorted maps:} Comparing different hashtable implementations (linear probing, closed addressing, etc.).
    \item \textbf{Sorted maps:} Comparing binary search trees, treaps, and splay trees.
    \item \textbf{Sorting routines:} Comparing Quicksort, Mergesort, and Heapsort against the standard library's \texttt{qsort}.
\end{itemize}

\section{Measuring memory usage}

In addition to time, we also monitored memory usage. We implemented a custom memory allocator wrapper (\texttt{libellul/src/memory.c}) that tracks the number of allocated blocks and the total reserved memory. This allows us to detect memory leaks and measure the memory overhead of different data structures.
